\chapter*{Ordliste}
\addcontentsline{toc}{chapter}{Ordliste}
\markboth{Ordliste}{}

Oppslagsverk over nøkkelomgrepa boka brukar. Definisjonane er bokas eigne: dei følgjer linja frå \emph{den dobbelte bokføringa} til \emph{Formlære}, og er difor ikkje alltid identiske med korleis faget sjølv brukar orda.

\begin{description}

\item[Affordans] Eigenskapen ved ei form som gjer bruksmoglegheitane hennar synlege: ein dørhank affordar drag, ein knapp affordar trykk (Norman). Boka les omgrepet som fagets fyrste tilløp til ein etterprøvbar standard, og som eitt trykk mellom mange, ikkje sjølve grunnen.

\item[Agent] Den som gjev form, forstått som ein navigatør i eit landskap med ein avgrensa horisont: han kan representere nokre aksar og er blind for andre. I rammeverket er designaren \emph{ein} agent blant mange (celle, materiale, marknad, organisasjon) og ikkje den einerådande.

\item[Brukarsentrering] Metodikken som flyttar blikket frå gjenstanden til mennesket framføre han. Boka viser at den same empatiske studien av brukaren som lova å lyfte faget, vart instrumentet for å kjenne henne godt nok til å fange henne.

\item[Den dobbelte bokføringa] Bokas grunnfigur: at designfaget fører \emph{inntektene} til ein vitskap (grader, professorat, autoritet) og \emph{utgiftene} til eit handverk (taus, lokal kompetanse) i den same boka. Ho fungerer berre så lenge ingen tel opp.

\item[Designerly ways of knowing] Cross sitt omgrep for ein eigen, ikkje-proposisjonell designkunnskap. Boka godtek skildringa, men peikar på at ho er bygd så ho alltid stemmer, og difor ikkje kan falsifiserast, og difor ikkje er eit fundament.

\item[Design thinking] Pakkinga av designpraksis (empati, idé, prototype, test) til ein allmenn problemløysingsmetode (IDEO, Brown). Boka les den enorme eksportsuksessen som beviset: metoden lét seg flytte til kven som helst nett fordi han var tom der eit fag skulle vere fullt.

\item[Falsifiserbarheit] Eigenskapen ved ein påstand at han kan visast gal (Popper). Boka brukar henne som skiljelina mellom eit fundament og ei vakker meining: det vrange problemet og \emph{designerly ways of knowing} kan ikkje ta feil; \emph{Formlære} kan, og det er deira einaste krav på å vere ein grunn.

\item[Form follows function] Den funksjonalistiske grunnsetninga (Sullivan). Boka reknar ho som tom: éi form følgjer ikkje av éin funksjon, fordi forma er underbestemd: ho er likevekta mellom mange samtidige trykk, og formelen skjuler dette i staden for å forklare det.

\item[Formgrammatikk] (\textit{shape grammar}) Eit sett generative reglar som, brukte om att på ei startform, produserer alle og berre dei formene eit system kan lage (Stiny og Gips, 1971). Ein etterprøvbar formalisme som gjorde «stil» om til ein regel som kan ta feil, og som faget valde mot.

\item[Formrom] (\textit{morphospace}) Rommet av moglege former, der kvar akse er ein målbar dimensjon. Forma er ikkje ein ting med eigenskapar, men eit punkt i dette rommet; dei tomme områda er like opplysande som dei fylte, for dei fortel kva seleksjonen forbyr.

\item[Forsking gjennom design] (\textit{research through design}) Ideen om at sjølve designpraksisen kan produsere kunnskap. Boka spør kva slik forsking etterlèt seg utan objekt og utan delt språk, og finn at ho oftast namngjev meir enn ho prøver.

\item[Funksjonalisme] Læra om at forma skal følgje av funksjonen. Boka handsamar han ikkje som ein teori om form, men som ein moral i forkledning, endå ein freistnad på å gje den tause formkunnskapen eit ord, som endar i ei overtyding.

\item[Før-paradigmatisk] Kuhns nemning på eit felt utan eit delt paradigme, der skular strir om grunnomgrepa og kunnskapen ikkje hopar seg opp. Boka plasserer designfaget her: ikkje umode på veg mot mognad, men strukturelt utan den presisjonen som skil dei paradigmatiske faga.

\item[Generativ modell] Ein modell som ikkje katalogiserer kva som finst, men genererer rommet av kva som er mogleg. I traktaten er forma skriven som ein slik modell: \(\mathrm{Form} = SG(\nabla_{C(A)} L(c,t)) \cap \lozenge\).

\item[Gute Form] Den tysk-sveitsiske etterkrigsideen om den objektivt gode forma (Bill, Rams). Boka les han som endå ein freistnad på å heve smak til standard utan eit fundament under, og difor like udefinert som den han skulle avløyse.

\item[Handverk] Kunnskapsmåten der tenking og laging bur i den same kroppen, ofte i den same rørsla; handa tenkjer (Sennett). Ikkje eit ufullstendig forstadium til design, men ein heil epistemisk struktur faget braut med og aldri kunne erstatte.

\item[Knowing how / knowing that] Ryles skilje mellom å kunne gjere noko og å vite at noko er tilfellet. Boka brukar det som drepande for draumen om formalisering: handverkets \emph{knowing how} er ikkje teori sett ut i livet, og kan ikkje vekslast inn i eit fags \emph{knowing that}.

\item[Kognitiv lyskjegle] (\(\lozenge\)) Den horisonten som avgrensar kva aksar ein agent kan representere (Fields og Levin). Det du ligg utanfor lyskjegla di, kan du verken spesifisere eller forsvare; her får \emph{det manglande nei-et} sitt presise uttrykk.

\item[Lauget] Institusjonen som bar formkunnskapen gjennom kroppar og lærlingordning, utan nedskriven teori. Boka framhevar at det visste \emph{korleis} utan å vite \emph{kvifor}, og at det aldri kravde meir; difor lét det seg ikkje arve saman med ein vitskaps autoritet.

\item[Merksemdsøkonomi] Den økonomiske logikken der menneskeleg merksemd er varen som vert seld til annonsørar (Wu, Williams). Boka koplar henne til \emph{sviktsignaturen}: systemet optimaliserer éin målbar akse og er blindt for resten.

\item[Det manglande nei-et] Bokas figur for at faget strukturelt ikkje kan avvise noko. Utan ein presis modell av kva ei form gjer (kva agentar ho rører, kva ho legg til side) manglar faget målestokken til å seie at \emph{dette} skal ikkje lagast. Det som ikkje kan målast, kan ikkje avvisast.

\item[Morfometri] Den kvantitative skildringa av form som posisjon og deformasjon i eit rom (frå Thompson). Boka brukar henne for å vise at form lét seg gjere til målbar storleik lenge før design vart fødd, i andre fag.

\item[Overvakingskapitalisme] Zuboffs namn på logikken som einsidig gjer menneskeleg erfaring om til åtferdsdata, og data om til prediksjonsprodukt selde i marknader for framtidig åtferd. Her er brukaren ikkje kunden, men råvara.

\item[Paradigme] Kuhns omgrep for det delte rammeverket av omgrep, metode og målestokkar som lèt eit fag hope opp kunnskap. Å bli paradigmatisk er overgangen designfaget har utsett i hundre år, og som \emph{Formlære} hevdar er mogleg.

\item[Planlagd forelding] Den medvitne konstruksjonen av kortliva form, så han må kjøpast på nytt. Boka les henne som funksjonalismens stille motsetnad i praksis: ein form som svikter langs aksar produsenten ikkje var bedd om å forsvare.

\item[Refleksjon-i-handling] Schöns omgrep for korleis den dugande praktikaren tenkjer medan han handlar. Boka aktar innsikta, men plasserer henne mellom dei vakre meiningane: ho forklarer kvifor presisjon er vanskeleg, i staden for å levere han.

\item[Risikohandverk vs. visseheitshandverk] Pyes skilje mellom arbeid der resultatet til kvar tid står på spel og avheng av utøvarens dømekraft (risiko), og arbeid der det er sikra på førehand av maskin eller mal (vissehet). Merksemda risikoen krev, er sjølv ein form for kunnskap.

\item[Seleksjonstrykk] Eit av dei samtidige, motstridande trykka som verkar på forma (ergonomi, materialøkonomi, kulturell meining, nedbryting, makt) og kvart dreg henne i si retning. Forma er likevekta mellom dei, ikkje følgja av eitt.

\item[Semantisk vending] Skiftet på 1980-talet til at ein gjenstands viktigaste eigenskap er kva han \emph{tyder}, ikkje kva han gjer (Krippendorff, Memphis). Boka viser at vendinga berre flytta udefinertheita frå funksjon til meining, og kalla flyttinga eit framsteg.

\item[Substrat-uavhengig] Eigenskapen ved agensmodellen at han ikkje bryr seg om kva agenten er laga av, berre kva rom han kan representere og kor langt horisonten rekk. Difor fell celle, handverkar, marknad og materiale inn under same skildring.

\item[Sviktsignatur] Bokas omgrep for at svikten til ei form er leseleg som signaturen av kva aksar agenten \emph{ikkje} representerte. Skaden er ikkje tilfeldig: han peikar tilbake på blindskapen som laga han.

\item[Systemorientert design] Retninga som vender designblikket mot heile system heller enn einskildobjekt. Boka rettar «kartet er ikkje ein teori»-kritikken mot henne: å teikne samanhengar er ikkje det same som å ha ein falsifiserbar modell av dei.

\item[Taus kunnskap] Den reelle, presise kompetansen som ber prestasjonen utan å la seg formulere: vi veit meir enn vi kan seie (Polanyi). Ekte kunnskap, men ikkje eit fundament; påkalla som forsvar blir han ei orsaking i staden for ein grunn.

\item[Tenestedesign] Faget for å forme og optimalisere tenester og system. Boka brukar det grensetilfellet der objektet er eit menneske (Palantir) til å vise at rammeverket inga grense har for korleis eit objekt kan handsamast, så lenge handsaminga kan kallast funksjonell.

\item[Tilpassingslandskap] (\textit{adaptive landscape}) Rommet der kvart punkt er ein mogleg tilstand og høgda er kor godt tilpassa han er til dei samtidige trykka (Wright, 1932). Underbestemminga faget opplevde som eit mysterium, er her sjølve strukturen: fleire toppar, fleire gode løysingar.

\item[Vrange problem] (\textit{wicked problems}) Rittel og Webbers omgrep for problem utan endeleg formulering og utan rett løysing. Boka hevdar at faget gjorde det om til eit alibi: sidan eit vrangt problem pr. definisjon ikkje kan løysast rett, kan ingen formgjeving motseiast, og presisjon vert unødvendig i staden for umogleg.

\end{description}
